\documentclass[12pt, a4paper]{ctexart}
\usepackage{amsmath, amssymb, amsthm}
\usepackage{fancyhdr}
\usepackage{geometry}
\usepackage{enumitem}
\usepackage{graphicx}
\usepackage{tikz}
\usepackage{pgfplots}
\pgfplotsset{compat=1.18}

\geometry{left=2cm, right=2cm, top=2.5cm, bottom=2.5cm}
\setCJKmainfont{Hiragino Sans GB}
\setCJKsansfont{Heiti SC}

\pagestyle{fancy}
\fancyhf{}
\fancyhead[L]{培优专题·不等式放缩技巧}
\fancyhead[R]{测试卷}
\fancyfoot[C]{\thepage}
\renewcommand{\headrulewidth}{0.4pt}

\newcommand{\score}[1]{\hfill\makebox[2cm][l]{(#1分)}}

\begin{document}

\begin{center}
\Large\textbf{培优专题测试卷·不等式放缩技巧}

\normalsize
考试时间：120分钟\quad 满分：150分
\end{center}

\vspace{0.5cm}

\noindent\textbf{一、单选题（本大题共8小题，每小题5分，共40分。在每小题给出的四个选项中，只有一项是符合题目要求的。）}

\vspace{0.3cm}

\begin{enumerate}
\item 已知 $a > b > 0$，则下列不等式恒成立的是\score{5}
\begin{enumerate}[label=\Alph*.]
\item $a + \dfrac{1}{b} > b + \dfrac{1}{a}$
\item $\dfrac{2a+b}{a+2b} > \dfrac{a}{b}$
\item $a^a b^b > a^b b^a$
\item $\ln a - \ln b > \dfrac{2(a-b)}{a+b}$
\end{enumerate}

\item 设 $x, y > 0$ 且 $x + y = 1$，则 $\dfrac{1}{x} + \dfrac{4}{y}$ 的最小值为\score{5}
\begin{enumerate}[label=\Alph*.]
\item $8$
\item $9$
\item $10$
\item $12$
\end{enumerate}

\item 已知 $n \in \mathbb{N}^*$，则 $\displaystyle\sum_{k=1}^{n} \frac{1}{k^2}$ 与 $2 - \dfrac{1}{n}$ 的大小关系是\score{5}
\begin{enumerate}[label=\Alph*.]
\item $\displaystyle\sum_{k=1}^{n} \frac{1}{k^2} < 2 - \frac{1}{n}$
\item $\displaystyle\sum_{k=1}^{n} \frac{1}{k^2} \leqslant 2 - \frac{1}{n}$
\item $\displaystyle\sum_{k=1}^{n} \frac{1}{k^2} > 2 - \frac{1}{n}$
\item 当 $n \geqslant 2$ 时，$\displaystyle\sum_{k=1}^{n} \frac{1}{k^2} < 2 - \frac{1}{n}$
\end{enumerate}

\item 若 $0 < x < \dfrac{\pi}{2}$，则 $x - \sin x$ 与 $\dfrac{x^3}{6}$ 的大小关系为\score{5}
\begin{enumerate}[label=\Alph*.]
\item $x - \sin x > \dfrac{x^3}{6}$
\item $x - \sin x < \dfrac{x^3}{6}$
\item $x - \sin x = \dfrac{x^3}{6}$
\item 大小关系不确定
\end{enumerate}

\item 已知 $a, b, c$ 为正实数且 $abc = 1$，则 $\dfrac{a}{ab + a + 1} + \dfrac{b}{bc + b + 1} + \dfrac{c}{ca + c + 1}$ 的值为\score{5}
\begin{enumerate}[label=\Alph*.]
\item $\dfrac{1}{3}$
\item $1$
\item $3$
\item 与 $a, b, c$ 的取值有关
\end{enumerate}

\item 设 $a = \ln 1.01$，$b = 0.01$，$c = \sqrt{1.01} - 1$，则\score{5}
\begin{enumerate}[label=\Alph*.]
\item $a < b < c$
\item $a < c < b$
\item $c < a < b$
\item $b < a < c$
\end{enumerate}

\item 已知正数 $a, b$ 满足 $a + b = 2$，则 $\dfrac{a^2 + 1}{a} \cdot \dfrac{b^2 + 1}{b}$ 的最小值为\score{5}
\begin{enumerate}[label=\Alph*.]
\item $2$
\item $\dfrac{5}{2}$
\item $\dfrac{9}{4}$
\item $3$
\end{enumerate}

\item 设 $f(x) = e^x - 1 - x - \dfrac{x^2}{2}$，当 $x > 0$ 时，$f(x)$ 的符号为\score{5}
\begin{enumerate}[label=\Alph*.]
\item $f(x) > 0$
\item $f(x) < 0$
\item $f(x) \geqslant 0$
\item $f(x) = 0$
\end{enumerate}

\end{enumerate}

\vspace{0.5cm}

\noindent\textbf{二、多选题（本大题共4小题，每小题5分，共20分。在每小题给出的选项中，有多项符合题目要求。全部选对得5分，部分选对得2分，有选错的得0分。）}

\vspace{0.3cm}

\begin{enumerate}[resume]
\item 设 $a, b > 0$，$a + b = 1$，则下列正确的有\score{5}
\begin{enumerate}[label=\Alph*.]
\item $ab \leqslant \dfrac{1}{4}$
\item $\sqrt{a} + \sqrt{b} \leqslant \sqrt{2}$
\item $\dfrac{1}{a} + \dfrac{1}{b} \geqslant 4$
\item $a^2 + b^2 \geqslant \dfrac{1}{2}$
\end{enumerate}

\item 已知 $n \in \mathbb{N}^*$ 且 $n \geqslant 2$，下列不等式成立的有\score{5}
\begin{enumerate}[label=\Alph*.]
\item $\left(1 + \dfrac{1}{n}\right)^n < 3$
\item $n! > \left(\dfrac{n}{2}\right)^{\frac{n}{2}}$
\item $\displaystyle\sum_{k=1}^{n} \frac{1}{\sqrt{k}} < 2\sqrt{n}$
\item $\left(1 + \dfrac{1}{n}\right)^{n+1} > \left(1 + \dfrac{1}{n+1}\right)^{n+2}$
\end{enumerate}

\item 关于函数 $f(x) = x\ln x$，下列说法正确的有\score{5}
\begin{enumerate}[label=\Alph*.]
\item $f(x)$ 在 $(0, +\infty)$ 上单调递增
\item 当 $x > 1$ 时，$x\ln x > x - 1$
\item 当 $0 < x < 1$ 时，$x\ln x > 1 - \dfrac{1}{x}$
\item 对任意 $a, b > 0$，$\dfrac{a\ln a - b\ln b}{a - b} > \ln\dfrac{a+b}{2}$（$a \neq b$）
\end{enumerate}

\item 已知正实数 $x, y, z$ 满足 $x + y + z = 1$，则下列正确的有\score{5}
\begin{enumerate}[label=\Alph*.]
\item $xyz \leqslant \dfrac{1}{27}$
\item $x^2 + y^2 + z^2 \geqslant \dfrac{1}{3}$
\item $\dfrac{1}{x} + \dfrac{1}{y} + \dfrac{1}{z} \geqslant 9$
\item $\sqrt{x} + \sqrt{y} + \sqrt{z} \leqslant \sqrt{3}$
\end{enumerate}

\end{enumerate}

\vspace{0.5cm}

\noindent\textbf{三、填空题（本大题共4小题，每小题5分，共20分。）}

\vspace{0.3cm}

\begin{enumerate}[resume]
\item 已知 $a, b > 0$ 且 $a + 4b = 4$，则 $\dfrac{1}{a} + \dfrac{1}{b}$ 的最小值为 \underline{\hspace{2cm}}。\score{5}

\item 设 $x > 0$，$y > 0$，$z > 0$，且 $x + y + z = 3$，则 $\dfrac{x^2}{y} + \dfrac{y^2}{z} + \dfrac{z^2}{x}$ 的最小值为 \underline{\hspace{2cm}}。\score{5}

\item 不等式 $\dfrac{x}{1+x^2} \leqslant \dfrac{1}{2}$ 对所有实数 $x$ 成立，等号成立的条件是 $x = $ \underline{\hspace{2cm}}。\score{5}

\item 已知 $a_n = \dfrac{n}{(n+1)!}$，则 $\displaystyle\sum_{n=1}^{20} a_n = $ \underline{\hspace{2cm}}（用阶乘表示）。\score{5}

\end{enumerate}

\vspace{0.5cm}

\noindent\textbf{四、解答题（本大题共6小题，共70分。解答应写出文字说明、证明过程或演算步骤。）}

\vspace{0.3cm}

\begin{enumerate}[resume]
\item （10分）已知 $a, b, c$ 为正实数，证明：
$$\frac{a^3}{b^2} + \frac{b^3}{c^2} + \frac{c^3}{a^2} \geqslant \frac{a^2}{b} + \frac{b^2}{c} + \frac{c^2}{a}$$

\item （12分）设 $n \in \mathbb{N}^*$，证明：
$$\frac{1}{2} \cdot \frac{3}{4} \cdot \frac{5}{6} \cdots \frac{2n-1}{2n} < \frac{1}{\sqrt{2n+1}}$$

\item （12分）已知函数 $f(x) = e^x - ax - 1$（$a \in \mathbb{R}$）。
\begin{enumerate}[label=(\arabic*)]
\item 当 $a = 1$ 时，证明：对任意 $x \geqslant 0$，$f(x) \geqslant \dfrac{x^2}{2}$；
\item 当 $a = 2$ 时，证明：对任意 $x \in \mathbb{R}$，$e^x \geqslant x^2 + 1$。
\end{enumerate}

\item （12分）已知正数数列 $\{a_n\}$ 满足 $a_1 = 1$，$a_{n+1} = a_n + \dfrac{1}{a_n}$（$n \in \mathbb{N}^*$）。
\begin{enumerate}[label=(\arabic*)]
\item 证明：$a_n^2 \geqslant 2n - 1$；
\item 证明：$\sqrt{2n} < a_n < \sqrt{2n} + \dfrac{\ln n}{2\sqrt{2n}}$（$n \geqslant 2$）。
\end{enumerate}

\item （12分）设 $a, b, c$ 为非负实数且 $a + b + c = 1$。
\begin{enumerate}[label=(\arabic*)]
\item 求 $a^2 b + b^2 c + c^2 a$ 的最大值；
\item 求 $\sqrt{a + \dfrac{(b-c)^2}{4}} + \sqrt{b + \dfrac{(c-a)^2}{4}} + \sqrt{c + \dfrac{(a-b)^2}{4}}$ 的最大值。
\end{enumerate}

\item （12分）已知 $f(x) = \ln(1+x) - \dfrac{x(1+\lambda x)}{1+x}$，其中 $\lambda \in \mathbb{R}$。
\begin{enumerate}[label=(\arabic*)]
\item 当 $\lambda = 0$ 时，证明：$f(x) \leqslant 0$ 对 $x > -1$ 恒成立；
\item 当 $\lambda \geqslant \dfrac{1}{2}$ 时，证明：$f(x) \leqslant 0$ 对 $x > -1$ 恒成立；
\item 设 $x_1, x_2, \ldots, x_n$ 为正实数且 $x_1 x_2 \cdots x_n = 1$，利用（2）的结论证明：$\displaystyle\sum_{i=1}^{n} \frac{1}{1+x_i} \leqslant 1 + \dfrac{n-1}{1 + \sqrt[n]{x_1 x_2 \cdots x_n}}$。\end{enumerate}

\end{enumerate}

\end{document}
